\notechapter{N-20251112}{Calculus Review I: Inequalities, Limits, Derivatives, and More}



\section{Cauchy--Schwarz and Minkowski}

The opening page recalls two inequalities:
\[
  \left(\sum_{i=1}^n a_i b_i
\right)^2
  \le
  \left(\sum_{i=1}^n a_i^2
\right)
  \left(\sum_{i=1}^n b_i^2
\right),
\]
and
\[
  \left(\sum_{i=1}^n(a_i+b_i)^2
\right)^{1/2}
  \le
  \left(\sum_{i=1}^n a_i^2
\right)^{1/2}
  +
  \left(\sum_{i=1}^n b_i^2
\right)^{1/2}.
\]
The first is the Cauchy--Schwarz inequality; the second is Minkowski's inequality, the triangle inequality in Euclidean norm.

The review function of these formulas is clear: they are often used to control error terms when proving limits.

\section{Neighborhood notation}

For $x_0\in\mathbb R$ and $\delta>0$, the $\delta$-neighborhood is
\[
  N_{\delta}(x_0)=\{x:|x-x_0|<\delta\}.
\]
The punctured neighborhood is
\[
  \mathring N_{\delta}(x_0)=\{x:0<|x-x_0|<\delta\}.
\]
The right and left neighborhoods are
\[
  N^+_{\delta}(x_0)=(x_0,x_0+\delta),\qquad
  N^-_{\delta}(x_0)=(x_0-\delta,x_0).
\]
For large variables, the notes use
\[
  U(\infty)=\{x:|x|>M\},\quad
  U(+\infty)=\{x:x>M\},\quad
  U(-\infty)=\{x:x<-M\}.
\]

\section{Trigonometric identities}

The first page lists a full table of trigonometric identities.  The basic ones are
\[
  \sin^2\alpha+\cos^2\alpha=1,\qquad
  1+	an^2\alpha=\sec^2\alpha,\qquad
  1+\cot^2\alpha=\csc^2\alpha.
\]
The addition formulas are
\[
  \sin(\alpha\pm\beta)=\sin\alpha\cos\beta\pm\cos\alpha\sin\beta,
\]
\[
  \cos(\alpha\pm\beta)=\cos\alpha\cos\beta\mp\sin\alpha\sin\beta,
\]
\[
  	an(\alpha\pm\beta)=
\frac{	an\alpha\pm	an\beta}{1\mp	an\alpha	an\beta}.
\]
The sum-to-product and product-to-sum formulas appear as well, for example
\[
  \sin\alpha+\sin\beta
  =2\sin
\frac{\alpha+\beta}{2}\cos
\frac{\alpha-\beta}{2},
\]
\[
  \sin\alpha\sin\beta
  =
\frac12[\cos(\alpha-\beta)-\cos(\alpha+\beta)].
\]

The page includes double-angle and power-reduction formulas:
\[
  \sin2\alpha=2\sin\alpha\cos\alpha,\qquad
  \cos2\alpha=2\cos^2\alpha-1=1-2\sin^2\alpha,
\]
\[
  \sin^2\alpha=
\frac12(1-\cos2\alpha),\qquad
  \cos^2\alpha=
\frac12(1+\cos2\alpha).
\]

\section{Hyperbolic functions}

The next page introduces functions generated by $e^x$:
\[
  \sinh x=
\frac12(e^x-e^{-x}),\qquad
  \cosh x=
\frac12(e^x+e^{-x}),
\]
\[
  	anh x=
\frac{e^x-e^{-x}}{e^x+e^{-x}},\qquad
  \coth x=
\frac{e^x+e^{-x}}{e^x-e^{-x}}.
\]
Their identities mirror trigonometric identities but with sign changes:
\[
  \cosh^2x-\sinh^2x=1,
\]
\[
  \sinh(x+y)=\sinh x\cosh y+\cosh x\sinh y,
\]
\[
  \cosh(x+y)=\cosh x\cosh y+\sinh x\sinh y.
\]
The note treats them as a useful parallel system rather than as isolated special functions.

\section{Sequence limits: epsilon--N language}

A sequence $\{x_n\}$ converges to $A$ if for every $\varepsilon>0$ there exists $N\in\mathbb N$ such that
\[
  n>N\quad\Longrightarrow\quad |x_n-A|<\varepsilon.
\]
We write
\[
  \lim_{n\to\infty}x_n=A.
\]
The note highlights the uniqueness theorem.  If a sequence had two limits $A$ and $B$, choose
\[
  \varepsilon=
\frac{|A-B|}{2}.
\]
For large $n$,
\[
  |A-B|\le |A-x_n|+|x_n-B|<2\varepsilon=|A-B|,
\]
a contradiction.

\section{Infinite limits and one-sided limits}

A sequence tends to $+\infty$ if for every $G>0$ there is $N$ such that
\[
  n>N\quad\Longrightarrow\quad x_n>G.
\]
For functions, the epsilon-delta definition is
\[
  \lim_{x\to x_0}f(x)=A
\]
if for every $\varepsilon>0$ there exists $\delta>0$ such that
\[
  0<|x-x_0|<\delta
  \quad\Longrightarrow\quad
  |f(x)-A|<\varepsilon.
\]
One-sided limits are defined by replacing the punctured neighborhood with a right or left punctured neighborhood.  A two-sided limit exists iff the two one-sided limits exist and are equal.

\section{Heine criterion for function limits}

The highlighted theorem says:
\[
  \lim_{x\to x_0}f(x)=A
\]
if and only if for every sequence $x_n\to x_0$ with $x_n\ne x_0$, one has
\[
  f(x_n)\to A.
\]
This is often used to prove that a limit does not exist.  For example,
\[
  \lim_{x\to0}\cos
\frac1x
\]
does not exist because along suitable sequences the values tend to different limits.  The scanned page uses the sequences
\[
  x_n=
\frac1{2n\pi},\qquad
  y_n=
\frac1{2n\pi+\pi/2}.
\]
Then $x_n,y_n\to0$, but $f(x_n)	o1$ and $f(y_n)	o0$.

\section{Infinitesimals and infinitesimal orders}

An infinitesimal is a variable whose limit is zero.  The note warns that a fixed small number such as $e^{-100}$ is not an infinitesimal; infinitesimal means a variable in a limiting process.

The pages record the usual rules:
\begin{itemize}
  \item a finite sum of infinitesimals is infinitesimal;
  \item a bounded variable times an infinitesimal is infinitesimal;
  \item a constant times an infinitesimal is infinitesimal;
  \item the product of two infinitesimals is infinitesimal;
  \item the product of infinitely many infinitesimals need not be handled by a naive rule.
\end{itemize}

The handwritten orange comment gives examples such as
\[
  1,
\frac12,
\frac13,\ldots,
\frac1n,\ldots
\]
and warns that infinitely many factors require separate analysis.  The problem is not philosophical; it is about whether one is taking a finite algebraic combination or a changing infinite product.

For two infinitesimals $\alpha,\beta$, one writes
\[
  \alpha=o(\beta)
  \quad\Longleftrightarrow\quad
  \lim
\frac{\alpha}{\beta}=0,
\]
and
\[
  \alpha\sim\beta
  \quad\Longleftrightarrow\quad
  \lim
\frac{\alpha}{\beta}=1.
\]
Equivalent infinitesimals may be substituted in products and quotients when the relevant limits exist.

\section{Squeeze theorem and monotone convergence}

The squeeze theorem is recorded as:
\[
  X\le Y\le Z,\qquad \lim X=\lim Z=A
  \quad\Longrightarrow\quad
  \lim Y=A.
\]
The monotone convergence theorem says that a monotone bounded sequence converges.  For example, an increasing sequence with an upper bound converges to its least upper bound; a decreasing sequence with a lower bound converges to its greatest lower bound.

The page gives the standard proof for the decreasing case: let $A$ be the infimum.  Since $A$ is the greatest lower bound, for every $\varepsilon>0$ there exists some $x_N<A+\varepsilon$, and monotonicity then gives $A\le x_n\le x_N<A+\varepsilon$ for $n>N$.

\section{Cauchy criteria}

The sequence Cauchy criterion is
\[
  \{x_n\}	ext{ converges}
  \quad\Longleftrightarrow\quad
\forall \varepsilon>0,\ \exists N,\ m,n>N\Rightarrow |x_m-x_n|<\varepsilon.
\]
For function limits, the analogous criterion is: $f(x)$ has a finite limit at $x_0$ iff for every $\varepsilon>0$ there is a punctured neighborhood such that
\[
  x_1,x_2\in \mathring N_\delta(x_0)
  \quad\Longrightarrow\quad
  |f(x_1)-f(x_2)|<\varepsilon.
\]
The note states that the proof of sufficiency uses completeness of the real numbers, so it is often skipped in elementary courses.

\section{Continuity}

A function is continuous at $x_0$ if
\[
  \lim_{x\to x_0}f(x)=f(x_0).
\]
Equivalent sequentially:
\[
  x_n\to x_0\quad\Longrightarrow\quad f(x_n)\to f(x_0).
\]
The pages list the usual operations: sums, products, quotients with nonzero denominator, and compositions of continuous functions are continuous where defined.

The note uses continuity to justify passing limits through elementary functions, for example
\[
  \lim \sin u_n=\sin(\lim u_n)
\]
when $u_n$ has a limit.

\section{Derivatives and differentials}

The later pages review derivative definitions:
\[
  f'(x_0)=\lim_{\Delta x\to0}
\frac{f(x_0+\Delta x)-f(x_0)}{\Delta x}.
\]
The differential form is
\[
  dy=f'(x)\,dx.
\]
The derivative rules are:
\[
  (u+v)'=u'+v',\qquad
  (uv)'=u'v+uv',\qquad
  \left(
rac uv
\right)'=
\frac{u'v-uv'}{v^2},
\]
\[
  (f\circ g)'(x)=f'(g(x))g'(x).
\]
The scanned pages also include logarithmic differentiation and implicit differentiation examples.  The practical method is to take logarithms when products, quotients, and powers are entangled, and to differentiate both sides when the function is defined implicitly.

\section{Taylor-type expansions}

The final pages include expansion examples and reminders about higher derivatives.  The standard Taylor form near $x=0$ is
\[
  f(x)=f(0)+f'(0)x+
\frac{f''(0)}{2!}x^2+\cdots.
\]
Frequently used expansions include
\[
  e^x=1+x+
\frac{x^2}{2!}+\cdots,
\]
\[
  \sin x=x-
\frac{x^3}{3!}+
\frac{x^5}{5!}-\cdots,
  \qquad
  \cos x=1-
\frac{x^2}{2!}+
\frac{x^4}{4!}-\cdots.
\]
The handwritten derivations use these expansions to compare infinitesimals and compute limits.

\section{Expanded synthesis and advanced context}

This note is a calculus foundation review.  The chain of ideas is:
\[
\begin{gathered}
  \text{neighborhoods}\to \text{epsilon definitions}\to \text{sequence criteria}\\
  \to \text{infinitesimals}\to \text{continuity}\to \text{derivatives and expansions}.
\end{gathered}
\]
The handwritten comments are mostly warnings: do not confuse small constants with infinitesimals; do not multiply infinitely many infinitesimals naively; do not prove nonexistence of a limit without constructing incompatible approaches; do not use equivalent infinitesimals outside their valid algebraic context.


\paragraph{Expanded synthesis.}
The common theme of this chapter is that an elementary limiting or computational rule becomes powerful only after one specifies the space in which the limiting process takes place. The earlier sections (`Cauchy--Schwarz and Minkowski`, `Neighborhood notation`, `Trigonometric identities`, `Hyperbolic functions`) may look like separate techniques, but they all ask the same question: which operations are stable under approximation? A derivative is stable first-order approximation,
\[
  f(a+h)=f(a)+L(h)+o(|h|),
\]
while an integral is stable accumulation with respect to a measure, and a convergent series is a limit in a chosen norm or topology.

The advanced bridge is functional-analytic. Uniform convergence is convergence in a sup norm; $L^p$ convergence is convergence in an integral norm; differentiability is the existence of a continuous linear approximation; many differential equations are operator equations $L[u]=f$. Theorems such as dominated convergence, the inverse function theorem, the projection theorem in Hilbert spaces, and semigroup existence results are refined answers to the same elementary concern: when may we pass from finite approximations to a genuine limiting object?

To use this viewpoint when rereading the original examples, identify the hidden stability condition. A Taylor expansion asks for control of the remainder. A change of variables asks for differentiability and Jacobian control. Termwise differentiation of a series asks for uniform convergence of derivative series. A differential-equation formula asks whether the operator is invertible under the imposed initial or boundary data. The elementary computation is the visible part; the advanced theory names the hypotheses that make it legitimate.
